
\subsection{诺维科夫环的哈密顿无关性}
\label{subsec:5-1-2}
\liftIndex{I347}

设 $\{\psi_t^{(0)}\}_{0\leq t\leq1}$
与 $\{\psi_t^{(1)}\}_{0\leq t\leq1}$
是两条哈密顿同痕，
分别用来移动 $L^{(0)}$ 与 $L^{(1)}$。
置
$\vec\psi=
(\{\psi_t^{(0)}\}_t,\{\psi_t^{(1)}\}_t)$。
对满足 $\ell(i)\in L^{(i)}$ 的道路
$\ell:[0,1]\to M$，
定义 $\ell^{\vec\psi}:[0,1]\to M$ 为
\[
\ell^{\vec\psi}(t)
=\bigl(\psi_t^{(1)}\circ\psi_{1-t}^{(0)}\bigr)
(\ell(t)).
\]
这诱导映射
\[
\begin{aligned}
I_{\vec\psi}
=I_{\psi^{(0)},\psi^{(1)}}:
\Omega(L^{(0)},L^{(1)})
\longrightarrow
\Omega\bigl(
\psi_1^{(0)}(L^{(0)}),
\psi_1^{(1)}(L^{(1)})\bigr).
\end{aligned}
\]
取目标道路空间中
包含
$\ell_0^{\vec\psi}=I_{\vec\psi}(\ell_0)$
的连通分支
\[
\Omega\bigl(
\psi_1^{(0)}(L^{(0)}),
\psi_1^{(1)}(L^{(1)});
\ell_0^{\vec\psi}\bigr),
\]
并以
\[
\Pi\bigl(
\psi_1^{(0)}(L^{(0)}),
\psi_1^{(1)}(L^{(1)});
\ell_0^{\vec\psi}\bigr)
\]
表示相应的覆盖变换群。
固定 $\ell_0^*\Lambda M$ 的一个截面
$\lambda^0$ 后，在
$(\ell_0^{\vec\psi})^*\Lambda M$
上考虑诱导截面
\[
\lambda^0_{\vec\psi}(t)
=\bigl(\psi_t^{(1)}\circ
       \psi_{1-t}^{(0)}\bigr)_*
\bigl(\lambda^0(t)\bigr).
\]

\begin{proposition}{5.1.7}\label{prop:5-1-7}
映射 $I_{\vec\psi}:\ell\mapsto\ell^{\vec\psi}$
诱导同构
\[
\begin{aligned}
I_{\vec\psi,*}:
\Pi(L^{(0)},L^{(1)};\ell_0)
\xrightarrow{\ \cong\ }
\Pi\bigl(
\psi_1^{(0)}(L^{(0)}),
\psi_1^{(1)}(L^{(1)});
\ell_0^{\vec\psi}\bigr),
\end{aligned}
\]
并满足下列关系\translatornote{原书实体页 279
把第二个等式的复合次序写反，
使原基道路上的 $\mu_{\lambda^0}$
无法作用于 $I_{\vec\psi,*}$ 的目标；
式 (5.1.13.2) 采用此处的型别相容次序。
见疑误 A-0259。}：
\[
E_{\ell_0^{\vec\psi}}\circ I_{\vec\psi,*}
=E_{\ell_0},
\qquad
\mu_{\lambda^0_{\vec\psi}}\circ I_{\vec\psi,*}
=\mu_{\lambda^0}.
\]
\end{proposition}

\begin{proof}
不失一般性，可假设
$\psi_\rho^{(0)}\equiv\mathrm{id}$；
此时把 $\vec\psi=(\mathrm{id},\psi)$
简记为 $\psi$。
下面的引理由
$\Omega(L^{(0)},L^{(1)};\ell_0)$
的定义立即推出。

\begin{lemma}{5.1.8}\label{lem:5-1-8}
设
$\gamma:S^1\to
\Omega(L^{(0)},L^{(1)};\ell_0)$
是环路，$\nu:S^1\times[0,1]\to M$
是由 $\gamma$ 诱导的相应映射。
则下列条件等价：
\par\noindent\textup{(5.1.9.1)}\quad
$\gamma$ 可提升为环路
$\widetilde\gamma:S^1\to
\widetilde\Omega(L^{(0)},L^{(1)};\ell_0)$。
\par\noindent\textup{(5.1.9.2)}\quad
$\displaystyle\int_{S^1\times[0,1]}\nu^*\omega=0$
且 $\mu([\gamma])=0$，
其中第 2.2.1 小节定义的马斯洛夫指标同态是
\[
\mu:\pi_1(\Omega(L^{(0)},L^{(1)}),\ell_0)
\longrightarrow\mathbb Z.
\]
\end{lemma}

对给定的哈密顿同痕
$\{\psi_\rho\}_{0\leq\rho\leq1}$，
考虑上述映射
\[
I_\psi:
\Omega(L^{(0)},L^{(1)})
\longrightarrow
\Omega(L^{(0)},\psi_1(L^{(1)})).
\]
为证明命题 5.1.7，
只须证明下述引理。

\begin{lemma}{5.1.10}\label{lem:5-1-10}
对给定的环路
$\gamma:S^1\to\Omega(L^{(0)},L^{(1)})$，
下列两个条件等价\translatornote{原书实体页 279
在 (5.1.11.2) 的未提升道路空间中
把第二端误写为 $\psi_1(L^{(0)})$；
同一分项的提升空间和上文 $I_\psi$
均要求 $\psi_1(L^{(1)})$。
见疑误 A-0260。}：
\par\noindent\textup{(5.1.11.1)}\quad
$\gamma$ 可提升为
$\widetilde\gamma:S^1\to
\widetilde\Omega(L^{(0)},L^{(1)})$。
\par\noindent\textup{(5.1.11.2)}\quad
$I_\psi(\gamma):
S^1\to\Omega(L^{(0)},\psi_1(L^{(1)}))$
可提升为环路
\[
\widetilde{I_\psi\circ\gamma}:
S^1\to
\widetilde\Omega(L^{(0)},\psi_1(L^{(1)})).
\]
\end{lemma}

\begin{proof}[引理 5.1.10 的证明]
假设 (5.1.11.1) 成立。
令 $\nu:S^1\times[0,1]\to M$
为 $\gamma$ 所对应的自然映射，
即 $\nu(\sigma,\rho)=\gamma(\sigma)(\rho)$。
显然，环路
$I_\psi(\gamma):
S^1\to\Omega(L^{(0)},\psi_1(L^{(1)}))$
对应的自然映射
$\nu_\psi:S^1\times[0,1]\to M$ 是
\begin{equation}
\nu_\psi(\sigma,\rho)
=\psi_\rho(\nu(\sigma,\rho))
=\psi_\rho(\gamma(\sigma)(\rho)).
\tag{5.1.12}\label{eq:5-1-12}
\end{equation}
由引理 5.1.8，
只须证明以下两个等式：
\begin{align}
\int\nu_\psi^*\omega
&=\int\nu^*\omega,
\tag{5.1.13.1}\label{eq:5-1-13-1}\\
\mu_{\lambda^0}([\gamma])
&=
\mu_{\lambda^0_{\{\mathrm{id},\psi\}}}
\bigl([I_\psi(\gamma)]\bigr).
\tag{5.1.13.2}\label{eq:5-1-13-2}
\end{align}

设 $H:[0,1]\times M\to\mathbb R$
是生成 $\psi$ 的哈密顿函数。
计算得
\[
\int\nu_\psi^*\omega
=\int_{S^1\times[0,1]}
\omega\left(
\frac{\partial\nu_\psi}{\partial\sigma},
\frac{\partial\nu_\psi}{\partial\rho}
\right)\,d\sigma\,d\rho.
\]
由式~\eqref{eq:5-1-12}
得到导数关系\translatornote{原书实体页 280
在本导数式及下式中均写 $X_H(\nu)$；
该向量位于 $T_\nu M$，
不能与位于 $T_{\nu_\psi}M$ 的
$d\psi_\rho(\partial_\rho\nu)$ 相加。
此处据式 (5.1.12) 恢复取值点。
见疑误 A-0264。}：
\[
\frac{\partial\nu_\psi}{\partial\sigma}
=d\psi_\rho
\left(\frac{\partial\nu}{\partial\sigma}\right),
\qquad
\frac{\partial\nu_\psi}{\partial\rho}
=X_H(\nu_\psi)
+d\psi_\rho
\left(\frac{\partial\nu}{\partial\rho}\right).
\]
其中 $d\psi_\rho$ 是 $\psi_\rho$ 的微分。
因而
\[
\begin{aligned}
\int\nu_\psi^*\omega
&=\int_{S^1\times[0,1]}
\omega\Bigl(
d\psi_\rho(\partial_\sigma\nu),\\[-1mm]
&\hspace{5em}
X_H(\nu_\psi)
+d\psi_\rho(\partial_\rho\nu)\Bigr)
\,d\sigma\,d\rho\\
&=\int_{S^1\times[0,1]}
\frac{\partial}{\partial\sigma}
H\bigl(\psi_\rho(\nu(\sigma,\rho))\bigr)
\,d\sigma\,d\rho
+\int\nu^*\omega\\
&=\int\nu^*\omega.
\end{aligned}
\]
这就证明了 (5.1.13.1)。

为证明 (5.1.13.2)，
考虑辛丛对
\[
\bigl(
\nu^*TM,\,
\nu_0^*TL^{(0)}
\sqcup\nu_1^*TL^{(1)}
\bigr)
\]
以及下列辛丛对\translatornote{原书实体页 280
第二个辛丛对的末端边界写成
$(\psi_1\circ\nu_1)^*TL^{(1)}$，
但 $\psi_1\circ\nu_1$ 取值于
$\psi_1(L^{(1)})$。
译文按丛对型别补出形变后的切丛。
见疑误 A-0261。}：
\[
\bigl(
\nu_\psi^*TM,\,
(\nu_\psi)_0^*TL^{(0)}
\sqcup
(\psi_1\circ\nu_1)^*
T(\psi_1(L^{(1)}))
\bigr).
\]
由重参数化映射族
\[
\nu_\psi^s(\sigma,\rho)
=\psi_{s\rho}(\nu(\sigma,\rho)),
\qquad s\in[0,1],
\]
直接可见这两个辛丛对同伦。
于是由辛丛对指标的同伦不变性，
\[
\mu(E_\nu,L_\nu)
=\mu(E_{\nu_\psi},L_{\nu_\psi}),
\]
从而得到 (5.1.13.2)。
\end{proof}

回到命题 5.1.7 的证明。
由引理 5.1.8 与刚证的引理 5.1.10\translatornote{原书实体页 280
此处只引用引理 5.1.8；
群映射下降还须使用引理 5.1.10
对可提升环路的保持性。
见疑误 A-0263。}，
上述映射 $I_\psi$
自然诱导覆盖变换群的映射
\[
I_\psi:
\Pi(L^{(0)},L^{(1)};\ell_0)
\longrightarrow
\Pi(L^{(0)},\psi_1(L^{(1)});
\ell_0^{\vec\psi}).
\]
相容关系
\[
E_{\ell_0^{\vec\psi}}\circ I_\psi
=E_{\ell_0}
\]
由 (5.1.13.1) 直接推出，
马斯洛夫指标的相容性则由
(5.1.13.2) 给出。
命题 5.1.7 得证。
\end{proof}

命题 5.1.7 蕴含下述结论。

\begin{corollary}{5.1.14}\label{cor:5-1-14}
存在同构
\[
\begin{aligned}
I_{\vec\psi,*}:
\Lambda(L^{(0)},L^{(1)};\ell_0)
\longrightarrow
\Lambda\bigl(
\psi_1^{(0)}(L^{(0)}),
\psi_1^{(1)}(L^{(1)});
\ell_0^{\vec\psi}\bigr),
\end{aligned}
\]
使得下式成立\translatornote{原书实体页 281
在本推论的相容性等式中
突然改用未定义的 $\vec\phi$；
本小节固定的数据一直是 $\vec\psi$。
见疑误 A-0262。}：
\[
I_{\psi_1^{(0)}(L^{(0)}),
   \psi_1^{(1)}(L^{(1)});
   \ell_0^{\vec\psi}}
\circ I_{\vec\psi,*}
=I_{L^{(0)},L^{(1)};\ell_0}.
\]
\end{corollary}
